% !TEX root = ../../../main.tex
% 所属：第一部分：单变量分析学 / 第5章：级数理论
% 节：5.5 无穷乘积
% 正文按原文件顺序保存；由 ../chapter5.tex 调用。
\section{无穷乘积}
\begin{definition}
	对数列 $\{p_n\}_{n=1}^{+\infty}$，且对所有 $n \ge 1$，$p_n \neq 0$，称
	\[
	p_1 p_2 \cdots p_n \cdots = \prod_{n=1}^{\infty} p_n
	\]
	为无穷乘积。定义其部分积数列为 $\{P_n\}$：
	\[
	\begin{aligned}
		P_1 &= p_1 \\
		P_2 &= p_1 p_2 \\
		P_3 &= p_1 p_2 p_3 \\
		&\vdots \\
		P_n &= p_1 p_2 \cdots p_n
	\end{aligned}
	\]
	
	若部分积数列 $\{P_n\}$ 收敛于某非零有限数 $P$，则称无穷乘积 $\displaystyle \prod_{n=1}^{\infty} p_n$ 收敛，并称 $P$ 为它的积，记为
	\[
	\prod_{n=1}^{\infty} p_n = P
	\]
	
	若 $\{P_n\}$ 发散，或 $\{P_n\}$ 收敛于 $0$，则称无穷乘积 $\displaystyle \prod_{n=1}^{\infty} p_n$ 发散。
\end{definition}
\begin{proposition}
	若无穷乘积 $\displaystyle \prod_{n=1}^{\infty} p_n$ 收敛，则
	\[
	\lim_{n \to \infty} p_n = 1
	\]
\end{proposition}
\begin{proof}
	设 $\displaystyle \prod_{n=1}^{\infty} p_n = P \neq 0$，那么
	\[
	\lim_{n \to \infty} p_n = \lim_{n \to \infty} \frac{P_n}{P_{n-1}} = \frac{P}{P} = 1
	\]
\end{proof}



从而，我们通常将 $p_n$ 写成 $1 + a_n$ 的形式，即
\[
\prod_{n=1}^{\infty} p_n = \prod_{n=1}^{\infty} (1 + a_n)
\]
由上述命题，若 $\displaystyle \prod_{n=1}^{\infty} p_n$ 收敛，则必有 $a_n \to 0$。

当 $n$ 充分大时，总有 $1 + a_n > 0$，而前有限项不影响无穷乘积的敛散性。因此，下面我们假设：
\begin{itemize}
	\item  对所有 $n$，$1 + a_n > 0$（即 $p_n > 0$）；
	
	\item $a_n \to 0$（即 $p_n \to 1$）。
\end{itemize}
在上述假设下，我们可以建立无穷乘积与无穷级数之间的关系。
\begin{proposition}\label{无穷乘积与无穷级数之间的关系}
	无穷乘积 $\displaystyle \prod_{n=1}^{\infty} p_n = \prod_{n=1}^{\infty} (1 + a_n)$ 收敛，当且仅当级数 $\displaystyle \sum_{n=1}^{\infty} \ln p_n = \sum_{n=1}^{\infty} \ln(1 + a_n)$ 收敛。
	
\end{proposition}
\begin{proof}
	设无穷乘积 $\prod p_n$ 的部分积为 $P_n$，级数 $\sum \ln p_n$ 的部分和为 $S_n$，则有
	\[
	\ln P_n = S_n \quad \Longleftrightarrow \quad P_n = e^{S_n}
	\]
	若 $S_n \to S$ 存在且有限，则 $\lim\limits_{n \to \infty} P_n = e^S$ 存在且有限，故 $\prod p_n$ 收敛。
	反之，若 $P_n \to P$ 存在、有限且大于 $0$，则 $\lim\limits_{n \to \infty} S_n = \ln P$ 存在且有限。
\end{proof}

\begin{corollary}\label{无穷乘积收敛}
	若对所有 $n \ge 1$，$a_n > 0$（或 $a_n < 0$），则无穷乘积 $\displaystyle \prod_{n=1}^{\infty} (1 + a_n)$ 收敛，当且仅当级数 $\displaystyle \sum_{n=1}^{\infty} a_n$ 收敛。
	
\end{corollary}
\begin{proof}
	考虑级数 $\sum \ln(1 + a_n)$ 与 $\sum a_n$。
	若对所有 $n \ge 1$，$a_n > 0$，则 $\ln(1 + a_n) > 0$。
	又因为
	\[
	\lim_{n \to \infty} \frac{\ln(1 + a_n)}{a_n} = 1 \quad (a_n \to 0)
	\]
	由正项级数的比较判别法，$\sum \ln(1 + a_n)$ 与 $\sum a_n$ 有相同的敛散性。
	而由\ref{无穷乘积与无穷级数之间的关系}，$\prod(1 + a_n)$ 与 $\sum \ln(1 + a_n)$ 有相同的敛散性，因此 $\prod(1 + a_n)$ 与 $\sum a_n$ 有相同的敛散性。
	
\end{proof}
\begin{corollary}
	对任意项级数 $\sum a_n$，若 $\sum a_n$ 收敛，则无穷乘积 $\displaystyle \prod_{n=1}^{\infty} (1 + a_n)$ 收敛，当且仅当 $\displaystyle \sum_{n=1}^{\infty} a_n^2$ 收敛。
	
\end{corollary}
\begin{proof}
	由不等式 $\ln(1 + a_n) \le a_n$，可得 $a_n - \ln(1 + a_n) \ge 0$。
	计算极限：
	\[
	\lim_{n \to \infty} \frac{a_n - \ln(1 + a_n)}{a_n^2} = \lim_{x \to 0} \frac{x - \ln(1 + x)}{x^2} = \frac{1}{2}
	\]
	故正项级数 $\sum (a_n - \ln(1 + a_n))$ 与 $\sum a_n^2$ 有相同的敛散性。
	
	若 $\prod(1 + a_n)$ 收敛，则由\ref{无穷乘积与无穷级数之间的关系}，$\sum \ln(1 + a_n)$ 收敛。又已知 $\sum a_n$ 收敛，故 $\sum (a_n - \ln(1 + a_n))$ 收敛，从而 $\sum a_n^2$ 收敛。反之同理可证。
\end{proof}
\begin{definition}
	若级数 $\sum \ln p_n = \sum \ln(1 + a_n)$ 绝对收敛，就称无穷乘积 $\displaystyle \prod_{n=1}^{\infty} p_n = \prod_{n=1}^{\infty} (1 + a_n)$ 绝对收敛。
\end{definition}
\begin{proposition}
	下面三个说法等价：
	\begin{enumerate}
		\item 无穷乘积 $\displaystyle \prod_{n=1}^{\infty} (1 + a_n)$ 绝对收敛；
		\item 无穷乘积 $\displaystyle \prod_{n=1}^{\infty} (1 + |a_n|)$ 收敛；
		\item 级数 $\displaystyle \sum_{n=1}^{\infty} |a_n|$ 收敛。
	\end{enumerate}
\end{proposition}
\begin{proof}
	由\ref{无穷乘积与无穷级数之间的关系}，$\prod(1 + |a_n|)$ 收敛当且仅当 $\sum \ln(1 + |a_n|)$ 收敛。又因为 $|a_n| \ge 0$，由推论 \ref{无穷乘积收敛}，$\sum \ln(1 + |a_n|)$ 收敛当且仅当 $\sum |a_n|$ 收敛，故 2 $\iff$ 3。
	
	下证 1 $\iff$ 3：
	根据定义，$\prod(1 + a_n)$ 绝对收敛 $\iff$ $\sum |\ln(1 + a_n)|$ 收敛。
	考虑极限：
	\[
	\lim_{n \to \infty} \frac{|\ln(1 + a_n)|}{|a_n|} = 1
	\]
	由正项级数的比较判别法，$\sum |\ln(1 + a_n)|$ 与 $\sum |a_n|$ 有相同的敛散性。因此 1 $\iff$ 3。
	
	综上，1、2、3 三个命题等价。
\end{proof}

下面介绍著名的 Wallis 公式与 Stirling 公式，涉及了一定的定积分，读者可以在读完定积分后再读，现在只需记下公式就足够了。
\begin{proposition}（Wallis 公式）
	设 \( p_n = 1 - \dfrac{1}{(2n)^2} \)，则
	\[
	\prod_{n=1}^{\infty} p_n = \frac{2}{\pi}
	\]
\end{proposition}
\begin{proof}
	\[
	\begin{aligned}
		P_n &= \prod_{k=1}^{n} \left(1 - \frac{1}{(2k)^2}\right)
		= \prod_{k=1}^{n} \frac{(2k-1)(2k+1)}{(2k)(2k)} \\
		&= \frac{1 \cdot 3 \cdot 3 \cdot \cdots \cdot (2n-1)(2n+1)}{2 \cdot 2 \cdot 4 \cdot 4 \cdot \cdots \cdot (2n)(2n)}
		= \frac{[(2n-1)!!]^2}{[(2n)!!]^2} \cdot (2n+1)
	\end{aligned}
	\]
	
	考虑定积分 \( I_n = \displaystyle\int_{0}^{\pi/2} \sin^n x \dd x \)。
	后面我们将知道：
	\[
	I_{2n} = \frac{(2n-1)!!}{(2n)!!} \cdot \frac{\pi}{2}, \quad I_{2n+1} = \frac{(2n)!!}{(2n+1)!!}
	\]
	因此
	\[
	\frac{\pi}{2} P_n = \frac{I_{2n}}{I_{2n+1}}
	\]
	
	由不等式 \( I_{2n+1} < I_{2n} < I_{2n-1} \)，得
	\[
	1 < \frac{I_{2n}}{I_{2n+1}} < \frac{I_{2n-1}}{I_{2n+1}}
	\]
	因为
	\[
	\lim_{n \to \infty} \frac{I_{2n-1}}{I_{2n+1}} = \lim_{n \to \infty} \frac{2n+1}{2n} = 1
	\]
	由夹逼性，
	\[
	\lim_{n \to \infty} P_n = \lim_{n \to \infty} \left( \frac{2}{\pi} \cdot \frac{I_{2n}}{I_{2n+1}} \right) = \frac{2}{\pi}
	\]
	故
	\[
	\prod_{n=1}^{\infty} \left(1 - \frac{1}{(2n)^2}\right) = \frac{2}{\pi}
	\]
\end{proof}


换一个形式，就是 Wallis 公式：
\[
\frac{\pi}{2} = \frac{2}{1} \cdot \frac{2}{3} \cdot \frac{4}{3} \cdot \frac{4}{5} \cdot \cdots \cdot \frac{2n}{2n-1} \cdot \frac{2n}{2n+1} \cdot \cdots
\]
\begin{proposition}
	（Stirling 公式）
	\[
	n! \sim \sqrt{2\pi n} \left( \frac{n}{e} \right)^n \quad (n \to +\infty)
	\]
	
\end{proposition}
\begin{proof}
	设 \( b_n = \dfrac{n! e^n}{n^{n+1/2}}, \ n=1,2,\dots \)，则
	\[
	\begin{aligned}
		\frac{b_n}{b_{n-1}} &= e\left(1 - \frac{1}{n}\right)^{n - \frac{1}{2}}
		= \exp\left[1 + \left(n - \frac{1}{2}\right)\ln\left(1 - \frac{1}{n}\right)\right] \\
		&= \exp\left( -\frac{1}{12n^2} + o\left(\frac{1}{n^2}\right) \right)
		= 1 - \frac{1}{12n^2} + o\left(\frac{1}{n^2}\right)
	\end{aligned}
	\]
	
	令 \( 1 + a_n = \dfrac{b_n}{b_{n-1}} \)，于是 \( \sum a_n \) 收敛。那么
	\[
	\prod_{n=2}^{\infty} (1 + a_n) = \prod_{n=2}^{\infty} \frac{b_n}{b_{n-1}}
	\]
	收敛，因此 \( \lim\limits_{n \to \infty} b_n = b_1 \displaystyle\prod_{n=2}^{\infty} \dfrac{b_n}{b_{n-1}} = A \neq 0 \)。
	
	利用 Wallis 公式：
	\[
	A = \lim_{n \to \infty} b_n = \lim_{n \to \infty} \frac{b_n^2}{b_{2n}}
	= \lim_{n \to \infty} \frac{(2n)!!}{(2n-1)!!} \cdot \sqrt{\frac{2}{n}} = \sqrt{2\pi}
	\]
	因此得到 Stirling 公式：
	\[
	n! \sim \sqrt{2\pi n} \left( \frac{n}{e} \right)^n
	\]
\end{proof}
\begin{example}
	证明
	\[
	\sin x=x\prod_{n=1}^{\infty}\left(1-\frac{x^2}{n^2\pi^2}\right).
	\]
\end{example}

\begin{proof}
	由三角函数的知识,我们知道
	\[
	\sin 3\varphi=\sin\varphi\left(3-4\sin^2\varphi\right),
	\]
	\[
	\sin 5\varphi=\sin\varphi\left(5-20\sin^2\varphi+16\sin^4\varphi\right).
	\]
	利用三角恒等式
	\[
	\sin(2k+1)\varphi
	=2(1-2\sin^2\varphi)\sin(2k-1)\varphi-\sin(2k-3)\varphi,
	\]
	以及应用数学归纳法,可以得到
	\[
	\sin(2n+1)\varphi=\sin\varphi\,P(\sin^2\varphi),
	\tag{*}
	\]
	其中 $P(u)$ 是 $n$ 次多项式,它的常数项为
	\[
	P(0)=\lim_{\varphi\to0}\frac{\sin(2n+1)\varphi}{\sin\varphi}=2n+1.
	\]
	由于
	\[
	\varphi=\frac{k\pi}{2n+1}\quad(k=1,2,\cdots,n)
	\]
	使 $(*)$ 式的左端成为 $0$,且
	\[
	\sin^2\frac{k\pi}{2n+1}\quad(k=1,2,\cdots,n)
	\]
	给出 $P(u)$ 的 $n$ 个不同的根,于是
	\[
	P(u)=(2n+1)\prod_{k=1}^{n}
	\left(1-\frac{u}{\sin^2\dfrac{k\pi}{2n+1}}\right).
	\]
	从而得到
	\[
	\sin(2n+1)\varphi
	=(2n+1)\sin\varphi\prod_{k=1}^{n}
	\left(1-\frac{\sin^2\varphi}{\sin^2\dfrac{k\pi}{2n+1}}\right).
	\]
	令 $x=(2n+1)\varphi$,代入后得到
	\[
	\sin x
	=(2n+1)\sin\frac{x}{2n+1}\prod_{k=1}^{n}
	\left(1-\frac{\sin^2\dfrac{x}{2n+1}}{\sin^2\dfrac{k\pi}{2n+1}}\right).
	\]
	因此,当 $m<n$ 时,成立
	\[
	\frac{\sin x}
	{(2n+1)\sin\dfrac{x}{2n+1}
		\displaystyle\prod_{k=1}^{m}
		\left(1-\frac{\sin^2\dfrac{x}{2n+1}}{\sin^2\dfrac{k\pi}{2n+1}}\right)}
	=
	\prod_{k=m+1}^{n}
	\left(1-\frac{\sin^2\dfrac{x}{2n+1}}{\sin^2\dfrac{k\pi}{2n+1}}\right).
	\tag{**}
	\]
	当 $n\to\infty$ 时,$(**)$ 式左端的极限为
	\[
	\frac{\sin x}{x\displaystyle\prod_{k=1}^{m}
		\left(1-\frac{x^2}{k^2\pi^2}\right)}.
	\]
	另一方面,取 $m$ 充分大,使得 $m>|x|/2$.由于
	\[
	\sin^2\frac{x}{2n+1}<\frac{x^2}{(2n+1)^2},
	\]
	并且当 $1\leq k\leq n$ 时,
	\[
	\sin^2\frac{k\pi}{2n+1}
	\geq \frac{4}{\pi^2}\frac{k^2\pi^2}{(2n+1)^2}
	=\frac{4k^2}{(2n+1)^2},
	\]
	故关于 $(**)$ 式的右端,有估计式
	\[
	\begin{aligned}
		1
		&>\prod_{k=m+1}^{n}
		\left(1-\frac{\sin^2\dfrac{x}{2n+1}}{\sin^2\dfrac{k\pi}{2n+1}}\right) \\
		&\geq \prod_{k=m+1}^{n}\left(1-\frac{x^2}{4k^2}\right)
		\geq \prod_{k=m+1}^{\infty}\left(1-\frac{x^2}{4k^2}\right).
	\end{aligned}
	\]
	令 $n\to\infty$,就得到
	\[
	1\geq
	\frac{\sin x}{x\displaystyle\prod_{k=1}^{m}
		\left(1-\frac{x^2}{k^2\pi^2}\right)}
	\geq
	\prod_{k=m+1}^{\infty}\left(1-\frac{x^2}{4k^2}\right).
	\]
	由级数 $\displaystyle\sum_{k=m+1}^{\infty}\frac{x^2}{4k^2}$ 的收敛性,可知无穷乘积
	\[
	\prod_{k=m+1}^{\infty}\left(1-\frac{x^2}{4k^2}\right)
	\]
	收敛,并且当 $m\to\infty$ 时,上述尾乘积趋于 $1$.根据夹逼性,得到
	\[
	\lim_{m\to\infty}
	\frac{\sin x}{x\displaystyle\prod_{k=1}^{m}
		\left(1-\frac{x^2}{k^2\pi^2}\right)}=1.
	\]
	于是
	\[
	\sin x=x\prod_{k=1}^{\infty}\left(1-\frac{x^2}{k^2\pi^2}\right).
	\]
	当 $x=0$ 或 $x=\pm n\pi$ 时,两端显然同为 $0$,故结论对所有 $x$ 成立.
\end{proof}






